\scanpage{111}
Données

a) anneau de base $k$

b) $k$-algèbres $A_1,A_2,A_3,A$

c) Sur chacune des $A_i$, structure quasi-quadratique sur $k$, i.e. on a
\diagram{\begin{tikzcd}A_i\arrow[r,shift left=2,"t_i"]\arrow[r,shift right=2,"n_i"']&k\end{tikzcd}}
\deleted{On pose}
satisfaisant les conditions $1^\circ$ à $5^\circ$ ci-dessus. On \unsure{notera} $a\longrightarrow\bar a$ l'involution canonique de $A_i$,
\[
 \bar a=t_i(a)-a
\]
\deleted{$\unsure{\bar a=\tau_i}$}
(NB on identifie \inserted{par \unsure{abus de} notation} $\lambda\in k$ avec son image dans les $A_i$, sans forcément supposer $k\longrightarrow A_i$ injectif\ldots)

d) Des Hom de $k$-algèbres
\diagram{\begin{tikzcd}[column sep=large]
 A_1\arrow[dr,"\alpha_1"']&A_2\arrow[d,"\alpha_2"]&A_3\arrow[dl,"\alpha_3"]\\
 &A&
\end{tikzcd}}
(pas néc. inj.), et des structures d'alg. quasi-quadratiques de $A$ sur les $A_i$
\deleted{\diagram{\begin{tikzcd}[column sep=large]
 &A\arrow[dl,shift right=2,"\unsure{T_1}"']\arrow[dl,shift left=2,"N_1"]\arrow[d,shift right=2,"\unsure{T_2}"']\arrow[d,shift left=2,"N_2"]\\
 A_1&\deleted{A_3}\inserted{\unsure{A_2}}&
\end{tikzcd}}}
\diagram{\begin{tikzcd}[column sep=large]
 &A\arrow[dl,shift right=2,"T_1"']\arrow[dl,shift left=2,"N_1"]\arrow[d,shift right=2,"T_2"']\arrow[d,shift left=2,"N_2"]\arrow[dr,shift right=2,"T_3"']\arrow[dr,shift left=2,"N_3"]&\\
 A_1&A_2&A_3
\end{tikzcd}}
, chacune satisfaisant donc au formulaire ci-dessus de $1^\circ$ à $5^\circ$. De plus, \uline{compatibilités}.
\deleted{
On pose
\[
 \sigma_i(x)=\unsure{T_i|\overline T_i}(x)-x
\]
\margin{$\sigma_i$ est un autom. involutif de $A$}
et on veut que
\[
 \boxed{\sigma_1\sigma_2=\sigma_3,\quad\sigma_2\sigma_3=\sigma_{\unsure{1|3}},\quad\sigma_3\sigma_1=\sigma_2}
\]
i.e. que les $\sigma_i$ sont des involutions qui \uline{commutent} entre elles, et $\boxed{\sigma_1\sigma_2\sigma_3=\operatorname{id}}$
}
\mspage{110}
